\section{The Five Diagram Families}
\label{sec:family-taxonomy}

The 22 Mermaid diagram types (plus our extensions) are organized into five families.
Each family shares a layout kernel, a visual vocabulary, and a set of Domain IR patterns.
This taxonomy drives architecture decisions: new types within a family reuse the existing
kernel and require primarily a new parser and minor IR extensions.

\subsection{Family Overview}

\begin{table}[h]
\centering\small
\caption{Five diagram families with constituent types and kernel}
\begin{tabularx}{\textwidth}{lXll}
\toprule
\textbf{Family} & \textbf{Diagram Types} & \textbf{Layout Kernel} & \textbf{Status} \\
\midrule
Node-link / Graph & flowchart, C4 (4 variants), architecture, block, requirement, gitGraph, sankey &
  Sugiyama (layered) & Flow built; rest planned \\
\addlinespace
UML / Software & sequence, class, state, erDiagram &
  Grammar-specific & Sequence built; rest Tier-1 \\
\addlinespace
Charts (data$\to$geometry) & pie, xychart, quadrant, radar &
  Grammar-of-graphics & New kernel (Tier-2) \\
\addlinespace
Timeline / Project & gantt, timeline, journey, kanban &
  Track-based & Timeline built; rest planned \\
\addlinespace
Tree / Hierarchy & mindmap, treemap &
  Buchheim--J\"unger--Leipert & Tree built; treemap planned \\
\bottomrule
\end{tabularx}
\label{tab:family-overview}
\end{table}

\subsection{Family 1: Node-Link / Graph}
\label{sec:family-nodelink}

The largest family. All members share the fundamental structure: \textbf{nodes} connected by
\textbf{edges}, laid out by a directed-graph algorithm.

\paragraph{Shared kernel.} The Sugiyama four-phase layered layout
(Section~\ref{sec:flow-grammar}): cycle removal, rank assignment, crossing minimization,
coordinate assignment. Variants (LR, TB, RL, BT) are direction parameters.

\paragraph{Type-specific extensions:}
\begin{description}
  \item[flowchart] Core type. Subgraphs, diverse node shapes, edge labels. \textbf{Built.}
  \item[C4 Context/Container/Component/Dynamic] Bounded contexts as styled node clusters;
        stereotyped relationships. Reuses flow kernel with C4-specific node rendering.
  \item[architecture] Icon-heavy nodes in a layered or grouped topology.
  \item[block] Block diagrams with nested containers---layout variation of flowchart.
  \item[requirement] Requirement nodes + satisfy/derive/refine edges.
  \item[gitGraph] Commit DAG with branch/merge topology. Specialized rank = commit order.
  \item[sankey] Weighted directed edges rendered as variable-width flows.
\end{description}

\subsection{Family 2: UML / Software (Tier-1 Priority)}
\label{sec:family-uml}

The differentiating family. These are the diagrams software teams use daily---and where
PlantUML's dated rendering creates the biggest opportunity.

\paragraph{Sequence diagrams.} Order-deterministic column layout. Participants, messages,
activations, fragments (alt/opt/loop/par/critical/break). \textbf{Built}
(Section~\ref{sec:sequence-grammar}).

\paragraph{Class diagrams.} Classes with attributes/methods, relationships
(inheritance, composition, aggregation, association, dependency), packages/namespaces.
Domain IR: classes[], relationships[]. Layout: hierarchical with inheritance edges
directing rank.

\paragraph{State diagrams.} States, transitions, composite states (nested), start/end
pseudo-states, fork/join bars, history nodes. Domain IR: states[], transitions[].
Layout: layered graph with composite-state nesting.

\paragraph{Entity-Relationship diagrams.} Entities with attributes, relationships with
cardinality (1, 0..1, 1..*, 0..*). Domain IR: entities[], relationships[].
Layout: force-directed or layered with relationship-label placement.

\subsection{Family 3: Charts (Data $\to$ Geometry)}
\label{sec:family-charts-overview}

A genuinely new grammar-of-graphics layer. Unlike the other families, charts map
\textbf{quantitative data} to \textbf{geometric marks} through scales and coordinate
systems. This is the Vega-Lite~\cite{vegalite2017} / Grammar of
Graphics~\cite{grammarofgraphics2005} lineage adapted to our architecture.

See Section~\ref{sec:chart-family} for the full specification.

\paragraph{Types:}
\begin{description}
  \item[pie] Proportional sectors from labeled values.
  \item[xychart] Bar and line charts on Cartesian axes.
  \item[quadrant] Four-quadrant scatter with labeled axes and positioned items.
  \item[radar] Multi-axis radial chart (spider/radar plot).
\end{description}

\subsection{Family 4: Timeline / Project}
\label{sec:family-timeline}

Track-based temporal layouts. All members share the concept of a time axis with items
positioned along it.

\paragraph{Shared kernel.} Track-based layout with configurable time axis, row allocation,
and activity/milestone placement (Section~\ref{sec:timeline-grammar}).

\paragraph{Types:}
\begin{description}
  \item[gantt] Tasks with durations, sections, dependencies (visual arrows), today marker.
        \textbf{Built} (maps to Timeline IR).
  \item[timeline] Sectioned event list on a time axis. \textbf{Built} (maps to Timeline IR).
  \item[journey] User-journey map with tasks scored by satisfaction level.
        Extension: satisfaction score $\to$ color/position mapping.
  \item[kanban] Column-based board (backlog/doing/done). Extension: columns as
        vertical tracks, cards as items.
\end{description}

\subsection{Family 5: Tree / Hierarchy}
\label{sec:family-tree}

Recursive parent-child structures laid out with tidy-tree algorithms.

\paragraph{Shared kernel.} Buchheim--J\"unger--Leipert O(n) tidy-tree layout
(Section~\ref{sec:tree-grammar}).

\paragraph{Types:}
\begin{description}
  \item[mindmap] Indentation-defined hierarchy with styled nodes. \textbf{Built.}
  \item[treemap] Area-proportional nested rectangles from weighted hierarchy.
        Extension: squarified treemap algorithm for the layout kernel.
\end{description}

\subsection{Coverage Roadmap Tiers}

\begin{description}
  \item[Tier 0 --- Wire existing grammars to Mermaid syntax.]
        flowchart, sequence, gantt, timeline, mindmap.
        Work: write parsers for Mermaid syntax $\to$ existing Domain IRs.
        \textbf{5 types. Kernel: already built.}

  \item[Tier 1 --- UML/software line.]
        class, state, erDiagram, C4 (4 variants).
        Work: new Domain IRs + layout extensions + Mermaid parsers.
        \textbf{7 types. High priority; key differentiator.}

  \item[Tier 2 --- Charts.]
        pie, xychart, quadrant, radar.
        Work: new chart grammar-of-graphics kernel + parsers.
        \textbf{4 types. New architecture required.}

  \item[Tier 3 --- Remaining types.]
        sankey, requirement, gitGraph, block, architecture, treemap, kanban, journey, packet.
        Work: mostly parser + IR variants on existing kernels.
        \textbf{9 types. Lower priority; incremental.}
\end{description}
